\section{An eigenvalue below the essential spectrum}
\label{sec:h-matter}

Proposition~\ref{prop:clique} is the sufficient condition from extra edges:
an induced \(K_n\), \(n\ge 10\), forces
\(m=-8-E_0\ge n-9\) on a finite perturbation of \(J_2\).
Proposition~\ref{prop:degree} says the pure \(8\)-regular host stops at
\(K_9\), whose trial value meets \(-8\) and does not go under it.
The same trial value clears the \(J_3\) floor \(-12\) only for \(n\ge 14\).
For \(10\le n\le 13\) it lies in the \(J_3\) essential spectrum.
For \(n\le 9\) the all-ones trial vector does not clear \(-8\).
In-band eigenvalues hybridize with the density \eqref{eq:dos}, which is the
width in Section~\ref{sec:h-splitdecay}.
The flat band at \(0\) is a kernel of the sheet-antisymmetric sector.
A packet supported there does not propagate in the quotient.
A quotient-moving mode sits in the dispersive band, and it remains an
eigenvalue outside that band only if some trial vector pushes it below the edge.
Lemma~\ref{lem:starsum} is the local form: an eigenstate has
\(\sum_{j\sim i}\mathcal{J}_{ij}=0\) at every vertex, and
\(\sum_{j\sim i}B_{ij}=-E|\psi_i|^2\).
A flat-band vector has vertex bond sum zero without each bond vanishing.
A state inside the band that fails the current sum can leak, with width
set by \eqref{eq:dos}.
